\section{Background}
\label{sec:background}

% Target: ~1 two-column page (~2 single-column equivalent)
% Condense report sections/background.tex significantly. Focus on what is
% essential for understanding the design and evaluation.

\subsection{Side-Channels and Covert Channels}
\label{sub:channels}

% [~0.5 column -- brief definitions]
%   - Side-channel: Unintended information leakage through observable
%     execution characteristics (timing, power, EM emanations).
%   - Covert channel: Intentional communication between colluding processes
%     (sender/Trojan and receiver/spy) that bypasses OS security boundaries.
%   - Distinction matters for threat modelling: side-channels are passive
%     observation; covert channels require active collusion.
%   - Running example: a naive Trojan that branches on a secret bit causes
%     observable timing difference in cache state; a Spectre-style gadget
%     uses speculative execution to encode secrets into the cache.

\subsection{OS-Level Timing Channels}
\label{sub:timing_channels}

% [~0.5 column -- three categories, one paragraph each]
%   1. Scheduler-based: Priority/preemptive scheduler makes task progress
%      observable. A high-priority task can encode information by choosing
%      to yield or spin; a low-priority spy measures its own wake latency.
%      This is the primary target of this work.
%   2. Shared-state (microarchitectural): Caches, TLBs, branch predictors,
%      and prefetchers retain state across context switches. A spy can probe
%      these structures after a victim runs.
%   3. Delay-based: A Trojan consumes a shared resource (e.g., memory bus,
%      disk I/O), and a spy measures increased latency.

\subsection{Spatial and Temporal Isolation}
\label{sub:isolation}

% [~0.5 column]
%   - Spatial isolation: Hardware enforcement of memory partitioning via
%     MMU/MPU. Well understood; FreeRTOS-MPU already provides per-task
%     memory protection.
%   - Temporal isolation: Resetting shared state to a known default between
%     security domains. This is the missing primitive that this work adds
%     to FreeRTOS.
%   - The two work together: spatial isolation prevents direct reads;
%     temporal isolation prevents indirect inference through timing.

\subsection{The \texttt{fence.t} Instruction}
\label{sub:fence_t}

% [~0.5 column]
%   - RISC-V extension: hardware instruction that drains on-core
%     microarchitectural state (L1 caches, TLBs, branch predictors, store
%     buffers) on in-order cores.
%   - fence.t.s variant for out-of-order CPUs with software support.
%   - Measured overhead in seL4: ~0.7\% on domain switches. Critical for
%     making temporal isolation practical.
%   - This work uses the QEMU emulation of fence.t (no real silicon
%     available at time of writing).

\subsection{FreeRTOS and FreeRTOS-MPU}
\label{sub:freertos}

% [~0.25 column]
%   - Widely deployed real-time OS for embedded systems. Default scheduler
%     is fixed-priority preemptive with round-robin among equal priorities.
%   - FreeRTOS-MPU variant provides per-task memory protection on
%     microcontrollers with a Memory Protection Unit. RISC-V MPU port used
%     as implementation platform.
%   - Existing MPU infrastructure (privileged/unprivileged tasks,
%     restricted task creation API) provides a natural foundation for
%     adding domain-level isolation.
